% Generated. Do not edit.
% Projection scriptural-date-location/section, version 1: the Scriptural Date and
% Location page, generated whole from research/chronology.toml, the Scripture
% chronology corpus and the source library by tools/tpt document-types refresh.
% document: liturgy/roman-rite/1962/propers/temporal/57-seventeenth-after-pentecost
\runninghead{Scriptural Date and Location}
\section{Scriptural Date and Location}
\label{sec:date-location}

{\small Each passage the Mass reads or sings from Scripture, in canonical order. The Context column names what is dated --- the event the passage narrates, a setting it remembers, the traditional attribution of its author, or the composition of its book --- and where it belongs; the Date column gives the answer of the Scripture chronology corpus for that subject, with every alternative it records. Texts of the Mass that are not drawn from Scripture have no row. The evidence for each date follows the table.\par}
\medskip
% Generated. Do not edit.
% Deterministic projection of research/chronology.toml.
% schema: 3
% document: liturgy/roman-rite/1962/propers/temporal/57-seventeenth-after-pentecost
% generated-by: tools/tpt proper-chronology annotations
\newcommand{\chronologyannotationclaim}[6]{#6}
\newcommand{\chronologyannotationcomparisonclaim}[8]{#8}
\newcommand{\chronologyannotationreach}[3]{}
\newcommand{\chronologyannotationgroup}[3]{#3}
\newcommand{\chronologyannotationcomparisongroup}[4]{#4}
\newcommand{\chronologyannotation}[1]{%
  \ifcsname triptychchronologyannotation@#1\endcsname
    \csname triptychchronologyannotation@#1\endcsname
  \else
    \PackageError{triptych}{No chronology annotation for #1}{}%
  \fi}
% comparison-dependency: src/claude/liturgy/roman-rite/1962/propers/temporal/57-seventeenth-after-pentecost/research/chronology-profile-comparisons.toml sha256=cf130edd51ee40f0ffa0e00307ea5df94bf8f8abcc31d08db70be52503ae6c64
% introit: Introit
\expandafter\def\csname triptychchronologyannotation@introit\endcsname{%
  \chronologyannotationgroup{traditional-attribution}{disputed}{\textbf{Traditional attribution} -- disputed: \chronologyannotationreach{Ps.118.1}{Ps.118}{true}\chronologyannotationreach{Ps.118.124}{Ps.118}{true}\chronologyannotationreach{Ps.118.137}{Ps.118}{true}\chronologyannotationclaim{israel.monarchy.david-traditional-era}{traditional-attribution}{catholic-traditional-v1}{disputed}{reigned from 1055 to 1015 B.C.}{David (reign in the usual chronology), B.C. 1055--1015}.}%
  \space \chronologyannotationgroup{historical-setting}{research-pending}{\textbf{Historical setting} -- Occasion date unresolved.}%
  \space \chronologyannotationgroup{composition}{preferred}{\textbf{Composition}: \chronologyannotationreach{Ps.118.1}{Ps}{true}\chronologyannotationreach{Ps.118.124}{Ps}{true}\chronologyannotationreach{Ps.118.137}{Ps}{true}\chronologyannotationclaim{critical.psalms.latest-composition-boundary}{composition}{catholic-critical-v1}{preferred}{before the Maccabean period, around 165 B.C.; no individual psalm can be dated securely}{Before c. 165 B.C.}}%
}
% epistle: Epistle
\expandafter\def\csname triptychchronologyannotation@epistle\endcsname{%
  \chronologyannotationgroup{composition}{disputed}{\textbf{Composition} -- disputed: \chronologyannotationreach{Eph.4.1}{Eph}{true}\chronologyannotationreach{Eph.4.2}{Eph}{true}\chronologyannotationreach{Eph.4.3}{Eph}{true}\chronologyannotationreach{Eph.4.4}{Eph}{true}\chronologyannotationreach{Eph.4.5}{Eph}{true}\chronologyannotationreach{Eph.4.6}{Eph}{true}\chronologyannotationclaim{composition.epistle-to-the-ephesians}{composition}{catholic-traditional-v1}{disputed}{a period between 58 and 63}{A.D. 58--63}; \chronologyannotationreach{Eph.4.1}{Eph}{true}\chronologyannotationreach{Eph.4.2}{Eph}{true}\chronologyannotationreach{Eph.4.3}{Eph}{true}\chronologyannotationreach{Eph.4.4}{Eph}{true}\chronologyannotationreach{Eph.4.5}{Eph}{true}\chronologyannotationreach{Eph.4.6}{Eph}{true}\chronologyannotationclaim{composition.epistle-to-the-ephesians}{composition}{catholic-traditional-v1}{disputed}{(Philemon; Colossians; Ephesians; Philippians), 61}{A.D. 61}.}%
}
% gradual: Gradual
\expandafter\def\csname triptychchronologyannotation@gradual\endcsname{%
  \chronologyannotationgroup{traditional-attribution}{disputed}{\textbf{Traditional attribution} -- disputed: \chronologyannotationreach{Ps.32.12}{Ps.32}{true}\chronologyannotationreach{Ps.32.6}{Ps.32}{true}\chronologyannotationclaim{israel.monarchy.david-traditional-era}{traditional-attribution}{catholic-traditional-v1}{disputed}{reigned from 1055 to 1015 B.C.}{David (reign in the usual chronology), B.C. 1055--1015}.}%
  \space \chronologyannotationgroup{historical-setting}{research-pending}{\textbf{Historical setting} -- Occasion date unresolved.}%
  \space \chronologyannotationgroup{composition}{preferred}{\textbf{Composition}: \chronologyannotationreach{Ps.32.12}{Ps}{true}\chronologyannotationreach{Ps.32.6}{Ps}{true}\chronologyannotationclaim{critical.psalms.latest-composition-boundary}{composition}{catholic-critical-v1}{preferred}{before the Maccabean period, around 165 B.C.; no individual psalm can be dated securely}{Before c. 165 B.C.}}%
}
% alleluia: Alleluia
\expandafter\def\csname triptychchronologyannotation@alleluia\endcsname{%
  \chronologyannotationgroup{historical-setting}{research-pending}{\textbf{Historical setting} -- Occasion date unresolved.}%
  \space \chronologyannotationgroup{composition}{preferred}{\textbf{Composition}: \chronologyannotationreach{Ps.101.2}{Ps}{true}\chronologyannotationclaim{critical.psalms.latest-composition-boundary}{composition}{catholic-critical-v1}{preferred}{before the Maccabean period, around 165 B.C.; no individual psalm can be dated securely}{Before c. 165 B.C.}}%
}
% gospel: Gospel
\expandafter\def\csname triptychchronologyannotation@gospel\endcsname{%
  \chronologyannotationgroup{narrated-event}{research-pending}{\textbf{Event} -- Narrated event date unresolved.}%
  \space \chronologyannotationgroup{composition}{disputed}{\textbf{Composition} -- disputed: \chronologyannotationreach{Matt.22.34}{Matt}{true}\chronologyannotationreach{Matt.22.35}{Matt}{true}\chronologyannotationreach{Matt.22.36}{Matt}{true}\chronologyannotationreach{Matt.22.37}{Matt}{true}\chronologyannotationreach{Matt.22.38}{Matt}{true}\chronologyannotationreach{Matt.22.39}{Matt}{true}\chronologyannotationreach{Matt.22.40}{Matt}{true}\chronologyannotationreach{Matt.22.41}{Matt}{true}\chronologyannotationreach{Matt.22.42}{Matt}{true}\chronologyannotationreach{Matt.22.43}{Matt}{true}\chronologyannotationreach{Matt.22.44}{Matt}{true}\chronologyannotationreach{Matt.22.45}{Matt}{true}\chronologyannotationreach{Matt.22.46}{Matt}{true}\chronologyannotationclaim{composition.gospel-of-matthew}{composition}{catholic-traditional-v1}{disputed}{about A.D. 38-45}{c. A.D. 38--45}; \chronologyannotationreach{Matt.22.34}{Matt}{true}\chronologyannotationreach{Matt.22.35}{Matt}{true}\chronologyannotationreach{Matt.22.36}{Matt}{true}\chronologyannotationreach{Matt.22.37}{Matt}{true}\chronologyannotationreach{Matt.22.38}{Matt}{true}\chronologyannotationreach{Matt.22.39}{Matt}{true}\chronologyannotationreach{Matt.22.40}{Matt}{true}\chronologyannotationreach{Matt.22.41}{Matt}{true}\chronologyannotationreach{Matt.22.42}{Matt}{true}\chronologyannotationreach{Matt.22.43}{Matt}{true}\chronologyannotationreach{Matt.22.44}{Matt}{true}\chronologyannotationreach{Matt.22.45}{Matt}{true}\chronologyannotationreach{Matt.22.46}{Matt}{true}\chronologyannotationclaim{composition.gospel-of-matthew}{composition}{catholic-traditional-v1}{disputed}{about the year 40-42}{c. A.D. 40--42}; \chronologyannotationreach{Matt.22.34}{Matt}{true}\chronologyannotationreach{Matt.22.35}{Matt}{true}\chronologyannotationreach{Matt.22.36}{Matt}{true}\chronologyannotationreach{Matt.22.37}{Matt}{true}\chronologyannotationreach{Matt.22.38}{Matt}{true}\chronologyannotationreach{Matt.22.39}{Matt}{true}\chronologyannotationreach{Matt.22.40}{Matt}{true}\chronologyannotationreach{Matt.22.41}{Matt}{true}\chronologyannotationreach{Matt.22.42}{Matt}{true}\chronologyannotationreach{Matt.22.43}{Matt}{true}\chronologyannotationreach{Matt.22.44}{Matt}{true}\chronologyannotationreach{Matt.22.45}{Matt}{true}\chronologyannotationreach{Matt.22.46}{Matt}{true}\chronologyannotationclaim{composition.gospel-of-matthew}{composition}{catholic-traditional-v1}{disputed}{the years 40-45}{A.D. 40--45}; \chronologyannotationreach{Matt.22.34}{Matt}{true}\chronologyannotationreach{Matt.22.35}{Matt}{true}\chronologyannotationreach{Matt.22.36}{Matt}{true}\chronologyannotationreach{Matt.22.37}{Matt}{true}\chronologyannotationreach{Matt.22.38}{Matt}{true}\chronologyannotationreach{Matt.22.39}{Matt}{true}\chronologyannotationreach{Matt.22.40}{Matt}{true}\chronologyannotationreach{Matt.22.41}{Matt}{true}\chronologyannotationreach{Matt.22.42}{Matt}{true}\chronologyannotationreach{Matt.22.43}{Matt}{true}\chronologyannotationreach{Matt.22.44}{Matt}{true}\chronologyannotationreach{Matt.22.45}{Matt}{true}\chronologyannotationreach{Matt.22.46}{Matt}{true}\chronologyannotationclaim{composition.gospel-of-matthew}{composition}{catholic-traditional-v1}{disputed}{about the year 60-68}{c. A.D. 60--68}; \chronologyannotationreach{Matt.22.34}{Matt}{true}\chronologyannotationreach{Matt.22.35}{Matt}{true}\chronologyannotationreach{Matt.22.36}{Matt}{true}\chronologyannotationreach{Matt.22.37}{Matt}{true}\chronologyannotationreach{Matt.22.38}{Matt}{true}\chronologyannotationreach{Matt.22.39}{Matt}{true}\chronologyannotationreach{Matt.22.40}{Matt}{true}\chronologyannotationreach{Matt.22.41}{Matt}{true}\chronologyannotationreach{Matt.22.42}{Matt}{true}\chronologyannotationreach{Matt.22.43}{Matt}{true}\chronologyannotationreach{Matt.22.44}{Matt}{true}\chronologyannotationreach{Matt.22.45}{Matt}{true}\chronologyannotationreach{Matt.22.46}{Matt}{true}\chronologyannotationclaim{composition.gospel-of-matthew}{composition}{catholic-traditional-v1}{disputed}{about the years 64-67}{c. A.D. 64--67}; \chronologyannotationreach{Matt.22.34}{Matt}{true}\chronologyannotationreach{Matt.22.35}{Matt}{true}\chronologyannotationreach{Matt.22.36}{Matt}{true}\chronologyannotationreach{Matt.22.37}{Matt}{true}\chronologyannotationreach{Matt.22.38}{Matt}{true}\chronologyannotationreach{Matt.22.39}{Matt}{true}\chronologyannotationreach{Matt.22.40}{Matt}{true}\chronologyannotationreach{Matt.22.41}{Matt}{true}\chronologyannotationreach{Matt.22.42}{Matt}{true}\chronologyannotationreach{Matt.22.43}{Matt}{true}\chronologyannotationreach{Matt.22.44}{Matt}{true}\chronologyannotationreach{Matt.22.45}{Matt}{true}\chronologyannotationreach{Matt.22.46}{Matt}{true}\chronologyannotationclaim{composition.gospel-of-matthew}{composition}{catholic-traditional-v1}{disputed}{about the year 50}{c. A.D. 50}.}%
  \space \chronologyannotationcomparisongroup{catholic-critical-v1}{composition}{composition-only}{\textbf{Composition (Catholic critical chronology, for comparison)}: \chronologyannotationreach{Matt.22.34}{Matt}{true}\chronologyannotationreach{Matt.22.35}{Matt}{true}\chronologyannotationreach{Matt.22.36}{Matt}{true}\chronologyannotationreach{Matt.22.37}{Matt}{true}\chronologyannotationreach{Matt.22.38}{Matt}{true}\chronologyannotationreach{Matt.22.39}{Matt}{true}\chronologyannotationreach{Matt.22.40}{Matt}{true}\chronologyannotationreach{Matt.22.41}{Matt}{true}\chronologyannotationreach{Matt.22.42}{Matt}{true}\chronologyannotationreach{Matt.22.43}{Matt}{true}\chronologyannotationreach{Matt.22.44}{Matt}{true}\chronologyannotationreach{Matt.22.45}{Matt}{true}\chronologyannotationreach{Matt.22.46}{Matt}{true}\chronologyannotationcomparisonclaim{critical.gospel-of-matthew}{composition}{catholic-critical-v1}{catholic-critical-v1}{preferred}{passage.united-states-conference-of-catholic-bishops.new-american-bible-revised-edition.english-usccb-web-2026-09-21.matthew-introduction}{post-A.D. 70 date}{The Greek Gospel of Matthew in the NABRE introduction, post-A.D. 70 date}.}%
}
% offertory: Offertory
\expandafter\def\csname triptychchronologyannotation@offertory\endcsname{%
  \chronologyannotationgroup{historical-setting}{preferred}{\textbf{Historical setting}: \chronologyannotationreach{Dan.9.17}{Dan}{true}\chronologyannotationreach{Dan.9.18}{Dan}{true}\chronologyannotationreach{Dan.9.19}{Dan}{true}\chronologyannotationclaim{israel.exile.seventy-years}{historical-setting}{catholic-traditional-v1}{preferred}{seventy years}{The seventy years of servitude to Babylon, ``seventy years'' (duration)}.}%
  \space \chronologyannotationgroup{composition}{preferred}{\textbf{Composition}: \chronologyannotationreach{Dan.9.17}{Dan}{true}\chronologyannotationreach{Dan.9.18}{Dan}{true}\chronologyannotationreach{Dan.9.19}{Dan}{true}\chronologyannotationclaim{composition.book-of-daniel}{composition}{catholic-traditional-v1}{preferred}{admits 570-536 B.C. as its date of composition}{B.C. 570--536}.}%
}
% communion: Communion
\expandafter\def\csname triptychchronologyannotation@communion\endcsname{%
  \chronologyannotationgroup{historical-setting}{research-pending}{\textbf{Historical setting} -- Occasion date unresolved.}%
  \space \chronologyannotationgroup{composition}{preferred}{\textbf{Composition}: \chronologyannotationreach{Ps.75.12}{Ps}{true}\chronologyannotationreach{Ps.75.13}{Ps}{true}\chronologyannotationclaim{critical.psalms.latest-composition-boundary}{composition}{catholic-critical-v1}{preferred}{before the Maccabean period, around 165 B.C.; no individual psalm can be dated securely}{Before c. 165 B.C.}}%
}

% Generated. Do not edit.
% Complete orientation-v3 table; wording, dates and order belong to the corpus.
\begin{chronologysharedtable}
\chronologyorientationrow{row-77de1d8acfc9b25e}{}
\chronologyorientationrow{row-051b97abaed87a71}{}
\chronologyorientationrow{row-3ea944ec9b52cbd0}{}
\midrule
\chronologyorientationrow{row-d53f96f04b0dbcba}{}
\chronologyorientationrow{row-51ede1dd606c99ee}{}
\midrule
\chronologyorientationrow{row-82bdd4fc18314432}{}
\chronologyorientationrow{row-3b5e6add4e39128f}{}
\midrule
\chronologyorientationrow{row-edbdd5308e9bef84}{}
\chronologyorientationrow{row-d43783bebe5506ad}{}
\chronologyorientationrow{row-904cdb00d172deab}{}
\midrule
\chronologyorientationrow{row-a6e80464fa57c53f}{}
\chronologyorientationrow{row-d4525052d5af80ed}{}
\midrule
\chronologyorientationrow{row-a13ab0733330eb08}{}
\chronologyorientationrow{row-5cfba241d7e648e3}{}
\chronologyorientationrow{row-bca6b7cfb8317ef6}{}
\midrule
\chronologyorientationrow{row-71a9502a065fbc47}{}
\end{chronologysharedtable}
{\footnotesize Psalm numbering: Vg. = Vulgate; Heb. = Hebrew.\par}


\subsection{Evidence for each date}
\label{sec:date-location-evidence}
{\small For each passage, as the Scriptural Date and Location table names it: what is dated, how the corpus weighs each answer, the source's own words, and the sources cited. A subject the corpus does not date is reported as such and given no date.\par}
{\small
\par\smallskip\noindent\textit{Gradual: Ps. 32:12 (Vg.; Heb. 33), Ps. 32:6 (Vg.; Heb. 33)}\par\nobreak
\begin{itemize}[leftmargin=1.2em,itemsep=1pt,topsep=1pt,parsep=0pt]
\item \textbf{Traditional attribution} (disputed). David's reign in the usual chronology: B.C.~1055--1015 (Corbett also reports writers dating the reign 30--50 years later) (``reigned from 1055 to 1015 B.C.''; The Catholic Encyclopedia, Volume 4 (New York, 1908), article 'David, King,' printed p. 642, opening paragraph). Place: Hebron, then Jerusalem (The Holy Bible, Douay-Rheims (Challoner revision), Project Gutenberg eBook 1581: 2 Kings 5:5).
\item \textbf{Historical setting}. No date: the chronology corpus has not yet dated this (research pending).
\item \textbf{Composition} (preferred). Psalms: Before c.~165 B.C. (``before the Maccabean period, around 165 B.C.; no individual psalm can be dated securely''; Catholic critical chronology; New American Bible, Revised Edition --- USCCB web state (2026), Psalms introduction).
\end{itemize}
\par\smallskip\noindent\textit{Communion: Ps. 75:12--13 (Vg.; Heb. 76)}\par\nobreak
\begin{itemize}[leftmargin=1.2em,itemsep=1pt,topsep=1pt,parsep=0pt]
\item \textbf{Historical setting}. As under Gradual: Ps. 32:12 (Vg.; Heb. 33), Ps. 32:6 (Vg.; Heb. 33).
\item \textbf{Composition}. As under Gradual: Ps. 32:12 (Vg.; Heb. 33), Ps. 32:6 (Vg.; Heb. 33).
\end{itemize}
\par\smallskip\noindent\textit{Alleluia: Ps. 101:2 (Vg.; Heb. 102)}\par\nobreak
\begin{itemize}[leftmargin=1.2em,itemsep=1pt,topsep=1pt,parsep=0pt]
\item \textbf{Historical setting}. As under Gradual: Ps. 32:12 (Vg.; Heb. 33), Ps. 32:6 (Vg.; Heb. 33).
\item \textbf{Composition}. As under Gradual: Ps. 32:12 (Vg.; Heb. 33), Ps. 32:6 (Vg.; Heb. 33).
\end{itemize}
\par\smallskip\noindent\textit{Introit: Ps. 118:137 (Vg.; Heb. 119), Ps. 118:124 (Vg.; Heb. 119), Ps. 118:1 (Vg.; Heb. 119)}\par\nobreak
\begin{itemize}[leftmargin=1.2em,itemsep=1pt,topsep=1pt,parsep=0pt]
\item \textbf{Traditional attribution}. As under Gradual: Ps. 32:12 (Vg.; Heb. 33), Ps. 32:6 (Vg.; Heb. 33).
\item \textbf{Historical setting}. As under Gradual: Ps. 32:12 (Vg.; Heb. 33), Ps. 32:6 (Vg.; Heb. 33).
\item \textbf{Composition}. As under Gradual: Ps. 32:12 (Vg.; Heb. 33), Ps. 32:6 (Vg.; Heb. 33).
\end{itemize}
\par\smallskip\noindent\textit{Offertory: Dan. 9:17--19}\par\nobreak
\begin{itemize}[leftmargin=1.2em,itemsep=1pt,topsep=1pt,parsep=0pt]
\item \textbf{Historical setting} (preferred). The seventy years of servitude to Babylon: ``seventy years'' (duration) (The Holy Bible, Douay-Rheims (Challoner revision), Project Gutenberg eBook 1581: Jer. 25:11, 25:12). Place: Judah and Babylon (The Holy Bible, Douay-Rheims (Challoner revision), Project Gutenberg eBook 1581: Jer. 25:11, 25:12).
\item \textbf{Composition} (preferred). The Book of Daniel: B.C.~570--536 (``admits 570-536 B.C. as its date of composition''; The Catholic Encyclopedia, Volume 4 (New York, 1908)). Place: Babylonian setting (The Holy Bible, Douay-Rheims (Challoner revision), Project Gutenberg eBook 1581: Dan. 1:1, 1:6, 7:1).
\end{itemize}
\par\smallskip\noindent\textit{Gospel: Matt. 22:34--46}\par\nobreak
\begin{itemize}[leftmargin=1.2em,itemsep=1pt,topsep=1pt,parsep=0pt]
\item \textbf{Narrated event}. No date: the chronology corpus has not yet dated this (research pending).
\item \textbf{Composition} (disputed). The Gospel of St Matthew: c.~A.D.~38--45 (``about A.D. 38-45''; The Catholic Encyclopedia, Volume 10 (1911); The Catholic Encyclopedia, Volume 14 (New York, 1912)); c.~A.D.~40--42 (``about the year 40-42''; The Catholic Encyclopedia, Volume 10 (1911)); A.D.~40--45 (``the years 40-45''; The Catholic Encyclopedia, Volume 10 (1911)); c.~A.D.~60--68 (``about the year 60-68''; The Catholic Encyclopedia, Volume 10 (1911)); c.~A.D.~64--67 (``about the years 64-67''; The Catholic Encyclopedia, Volume 10 (1911)); c.~A.D.~50 (``about the year 50''; The Catholic Encyclopedia, Volume 14 (New York, 1912)). Place: Semitic original associated with Palestine (The Catholic Encyclopedia, Volume 10 (1911)).
\item \textbf{Composition, Catholic critical comparison} (preferred). The Greek Gospel of Matthew in the NABRE introduction: Post-A.D.~70 date (Catholic critical chronology; New American Bible, Revised Edition: USCCB web state (2026), Matthew introduction, final five prose paragraphs on authorship, sources, composition and place; narrative overview of chapters 19-25). Place: Antioch suggested (New American Bible, Revised Edition: USCCB web state (2026), Matthew introduction, final five prose paragraphs on authorship, sources, composition and place; narrative overview of chapters 19-25).
\end{itemize}
\par\smallskip\noindent\textit{Epistle: Eph. 4:1--6}\par\nobreak
\begin{itemize}[leftmargin=1.2em,itemsep=1pt,topsep=1pt,parsep=0pt]
\item \textbf{Composition} (disputed). The Epistle to the Ephesians: A.D.~58--63 (``a period between 58 and 63''; The Catholic Encyclopedia, Volume 5 (New York, 1909)); A.D.~61 (``(Philemon; Colossians; Ephesians; Philippians), 61''; The Catholic Encyclopedia, Volume 11 (New York, 1911)). Place: Rome or Caesarea; recipients in Asia Minor (The Catholic Encyclopedia, Volume 5 (New York, 1909)).
\end{itemize}
\par}
